\documentclass[twocolumn]{aastex631}
\begin{document}
\title{The reionization ionizing-photon-budget shortfall for star-forming galaxies is not robust to systematics at z\~{}5}
\author{NebulaMind Lab (autonomous pipeline)}
\affiliation{NebulaMind Lab, \url{https://lab.nebulamind.net}}
\begin{abstract}
Reionization ionizing-photon-budget reconciliation at z\~{}5: star-forming galaxies require f\_esc=0.025 (+0.025/-0.013) to close the budget (Madau-Dickinson SFRD, log xi\_ion=25.5+/-0.15, clumping C=2-5, JWST-SFRD tail), versus indirect-proxy-inferred f\_esc=0.062 (+0.108/-0.039) from LzLCS O32/beta calibrations. Median delta(required-inferred)=-0.034 dex-frac (16-84\%: -0.142 to +0.008); 22\% of the systematic MC shows a shortfall. Conclusion: the budget CLOSES within the systematic, and the sign holds under both O32 and beta calibrations. The result is bounded by the xi\_ion x clumping x proxy-calibration systematic, not by statistics. Generated autonomously from public data (jwst) via the NebulaMind Lab runner: a bounded, reproducible, descriptive study.
\end{abstract}
\section{Introduction}
Recent studies have highlighted a potential crisis in reconciling the ionizing photon budget during reionization, particularly with the advent of new observational data from advanced telescopes [Muñoz2024]. The challenge lies in understanding whether star-forming galaxies alone can account for the necessary ionizing photons to drive reionization. Previous efforts have explored various aspects of this problem, including the role of excursion set models [Park2022], galaxy ionizing photon budgets at z < 10 [Duncan2015], and the demands on ionizing sources during absorption-dominated reionization [Davies2021]. To address this issue, we rely on established analytic frameworks for cosmic reionization [Madau2017].
\section{Data and method}
In our analysis, we adopt a literature-anchored approach to calculate the reionization photon budget. We utilize the Madau \& Dickinson (2014) cosmic star formation rate density (SFRD) as our foundation and incorporate published values for xi\_ion and O32/beta f\_esc proxy calibrations from LzLCS studies [Chisholm+22, Flury+22; Simmonds+24]. Notably, we do not employ any survey catalog data or direct observational inputs from JWST, SDSS, or TNG in this work. Instead, our method focuses on reconciling systematic uncertainties within the existing literature to assess the ionizing photon budget.
\section{Result}
Our key finding is that star-forming galaxies can close the reionization photon budget at z\~{}5 if they exhibit an escape fraction of f\_esc = 0.025 (+0.025/-0.013). This result is derived from combining the Madau-Dickinson SFRD, log xi\_ion=25.5+/-0.15, and clumping factor C=2-5 with JWST-SFRD tail considerations. In contrast, indirect-proxy-inferred escape fractions from LzLCS O32/beta calibrations suggest a value of f\_esc = 0.062 (+0.108/-0.039). The median difference between the required and inferred values is -0.034 dex-frac (16-84\%: -0.142 to +0.008), with 22\% of systematic Monte Carlo simulations indicating a shortfall in the photon budget.
\begin{figure}\centering\includegraphics[width=\columnwidth]{result.png}\caption{The reionization ionizing-photon-budget shortfall for star-forming galaxies is not robust to systematics at z\~{}5}\end{figure}
\section{Caveats}
It is essential to acknowledge the limitations of our approach. Our analysis relies on automated, single-selection, uncalibrated measurements, which may introduce biases and uncertainties not fully captured by our systematic error estimates. Furthermore, we depend heavily on previously published calibrations for xi\_ion and f\_esc proxy relationships, which may not account for all relevant astrophysical processes or observational complexities. Additionally, the clumping factor C remains a significant source of uncertainty in our calculations, as it is challenging to constrain observationally. These factors underscore the need for further research and refined measurements to solidify our understanding of the reionization photon budget.
\section*{References}
\footnotesize
[Muoz2024] Muñoz et al. (2024). Reionization after JWST: a photon budget crisis?. bibcode 2024MNRAS.535L..37M \par [Park2022] Park et al. (2022). Calibrating excursion set reionization models to approximately conserve ionizing photons. bibcode 2022MNRAS.517..192P \par [Duncan2015] Duncan et al. (2015). Powering reionization: assessing the galaxy ionizing photon budget at z < 10. bibcode 2015MNRAS.451.2030D \par [Davies2021] Davies et al. (2021). The Predicament of Absorption-dominated Reionization: Increased Demands on Ionizing Sources. bibcode 2021ApJ...918L..35D \par [Madau2017] Madau (2017). Cosmic Reionization after Planck and before JWST: An Analytic Approach. bibcode 2017ApJ...851...50M

\end{document}
